\section{Engineering the search}
\label{sec:engineering}
Most of the 2011 method's advantage turns out to be recoverable by engineering the same search. Three
changes were tested, each first validated against exact counts.

\paragraph{Bit-parallel candidate rows.} The search of Section~\ref{sec:tree} tests each candidate row
against every earlier column, $O(n\cdot c)$ work to expand a node at column $c$. Row $r$ conflicts with
column $i$ exactly when the difference $r-p_i$ is already used at lag $c-i$, so the rows forbidden by
column $i$ are the used-difference mask of that lag, shifted into row coordinates. ORing these masks gives
every forbidden row at once, in $O(c)$ word operations per node. The tree is unchanged: node and
solution counts are identical at orders 10 to 14. At order 32 the rate rises from $3.33\times10^6$ to
$1.68\times10^7$ nodes per second on the same core, a factor of $5.0$.

\paragraph{Column order.} Four orders were compared, all with the same symmetry restriction, which applies
to whichever column is placed first.

\begin{table}[H]
\centering
\small
\begin{tabular}{@{}rrrrr@{}}
\toprule
order & left to right & outside in & middle out & fewest legal rows \\
\midrule
12 & $3.49\times10^6$ ($0.36$\,s) & $1.22\times10^7$ ($1.39$\,s) & $3.49\times10^6$ ($0.44$\,s) & $2.42\times10^6$ ($1.22$\,s) \\
13 & $1.83\times10^7$ ($2.19$\,s) & $7.81\times10^7$ ($8.86$\,s) & $1.79\times10^7$ ($1.95$\,s) & $1.22\times10^7$ ($7.58$\,s) \\
14 & $8.32\times10^7$ ($12.5$\,s) & $4.38\times10^8$ ($55.9$\,s) & $8.32\times10^7$ ($11.5$\,s) & $5.38\times10^7$ ($35.6$\,s) \\
\bottomrule
\end{tabular}
\caption{Nodes and time for four column orders. ``Fewest legal rows'' chooses at each node the unplaced
column with fewest legal rows and abandons a branch when any column has none.}
\label{tab:orders}
\end{table}

\noindent Any order that keeps the placed columns contiguous gives the same tree, because whether a block
of dots is Costas does not depend on where the block sits; at even orders the counts for left to right and
middle out are identical. Placing dots far apart, outside in, wastes the short-lag constraints, which prune
hardest, and gives a tree five times larger by order 14. Choosing the column with fewest legal rows shrinks
the tree increasingly, to $65\%$ of the baseline at order 14, and its growth per order is about $4.65$
against $4.84$. But each node costs about three times as much, so it is slower at every order measured.

Orders that place a column between two earlier ones need one more check than left to right. A new dot
forms two pairs at the same lag $L$, one on each side, and these collide exactly when its row is the
midpoint of the other two, $2r=p_{c-L}+p_{c+L}$: the Ladendorf condition of the 2009 search. A first
version of these searches omitted the check and accepted non-Costas arrays; the validation against exact
array counts caught it before any result was used.

\paragraph{Full symmetry pruning.} Using all eight symmetries of the square, the search keeps exactly one
array per class: the lexicographically smallest of its eight images. A prefix is cut when some image is
guaranteed to be smaller for every completion. An image entry is either known, or confined to an interval:
an unplaced row's column lies among the unplaced columns. The vertical flip compares known values; the four
symmetries involving transposition can often be decided early from the partly known inverse; the horizontal
flip and the half-turn are decided only near the leaves. The pruning is exact: at every order from 4 to
13 the search outputs precisely the class representatives found by enumerating every Costas array and
grouping them by symmetry, including 990 classes among 7{,}852 arrays at order 12 and 1{,}616 among 12{,}828
at order 13. On the pruned tree the estimator is within $0.3\%$ of the exact counts at orders 13 to 15.

\begin{table}[H]
\centering
\small
\begin{tabular}{@{}rrrr@{}}
\toprule
order & one symmetry & all eight & fewer nodes \\
\midrule
12 & $3{,}488{,}680$ & $1{,}544{,}404$ & $2.26\times$ \\
13 & $18{,}290{,}043$ & $7{,}509{,}524$ & $2.44\times$ \\
14 & $83{,}176{,}084$ & $37{,}132{,}082$ & $2.24\times$ \\
15 & $435{,}090{,}280$ & $182{,}409{,}077$ & $2.39\times$ \\
\bottomrule
\end{tabular}
\caption{Nodes with one and with all eight symmetries, in the bit-parallel search.}
\label{tab:symmetry}
\end{table}

\noindent At order 32 three independent estimates put the pruned tree at $7.22\times10^{19}$ nodes, $2.13$
times smaller, while the extra comparisons lower the rate to $1.42\times10^7$ nodes per second. The search is
about $1.8$ times faster overall: less than the theoretical four, because two of the symmetries can be
decided only near the leaves.

\paragraph{The revised cost.} With the fully loaded slowdown of Section~\ref{sec:cluster}:

\begin{table}[H]
\centering
\small
\begin{tabular}{@{}lrrr@{}}
\toprule
search & core-hours & core-years & KSU \\
\midrule
as in Section~\ref{sec:cost} & $1.36\times10^{10}$ & $1.55\times10^6$ & $2.7\times10^7$ \\
bit-parallel candidate rows & $3.0\times10^{9}$ & $3.4\times10^5$ & $6.0\times10^6$ \\
and full symmetry pruning & $1.7\times10^{9}$ & $1.9\times10^5$ & $3.3\times10^6$ \\
\bottomrule
\end{tabular}
\caption{The cost of the order-32 search, before and after engineering.}
\label{tab:revised}
\end{table}

\noindent The engineered search is about eight times cheaper, but still about 18 years on $10^4$ cores. At
order 29 it would need about 1{,}650 core-years against the 2011 enumeration's 366.55 CPU-years, so a factor
of four to five remains, more if 2010 hardware was slower per core. Whether that factor lies in node rate or
in a smaller tree is not known. Reductions in the growth per order compound: the fewest-legal-rows order
lowers it from $4.84$ to about $4.65$, which over the orders up to 32 is worth about three times, if its
per-node cost could be brought down to the baseline's.

